\documentclass[10pt]{article}
\usepackage[margin=20mm]{geometry}
\usepackage{amsmath,amssymb,mathtools,fontspec,unicode-math,xurl,hyperref}
\IfFontExistsTF{DejaVu Serif}{\setmainfont{DejaVu Serif}}{\setmainfont{Latin Modern Roman}}
\setmathfont{latinmodern-math.otf}
\setlength{\parindent}{0pt}\setlength{\parskip}{6pt}\sloppy
\title{Technical Verification of Semiclassical Inequalities for Dirichlet and Neumann Laplacians on Convex Domains}
\author{AI-assisted technical review}
\date{}
\begin{document}
\maketitle
Abstract: Lemma 3.3 has a false quantitative conclusion: its coefficient omits the factor \((\gamma')^{\gamma'}/\gamma^\gamma\). An explicit square, with only its first Dirichlet eigenvalue contributing, satisfies every hypothesis and violates the asserted bound by the exact factor \(2^{1/8}\). The same coefficient occurs in the acquired journal version and arXiv v2. The corrected estimate preserves the local exponent improvement used in the qualitative extrapolation argument. In particular, this finding does not disprove the principal semiclassical inequalities below exponent one, the critical-exponent conclusions, or the shape-optimization theorems.\par
\section*{Identity, source and scope}
The reviewed paper is by Rupert L. Frank and Simon Larson, Communications on Pure and Applied Mathematics 79 (2026), 762–822, DOI 10.1002/cpa.70019 (https://doi.org/10.1002/cpa.70019) [FL]. The acquired 62-page PDF consists of a Chalmers cover followed by the 61-page version of record. Thus physical PDF page \(p\geq2\) corresponds to printed journal page \(p+760\). The cover and journal first page identify the version and authors. Its placement in the assigned annual batch is retained.\par
Root scientific reading covered physical pages 1–34 and 39–42. The detailed audit concentrated on the definitions and principal statements, the improved inequalities in Section 2, the full extrapolation argument in Section 3, the critical-exponent mechanism in Section 4, the product reduction in Section 5.1, and the Dirichlet dependency in the proof of Theorem 1.5. The remaining optimization and collapsing-domain proofs were not independently reconstructed. This report makes no certification of those unread parts or of every cited external result.\par
For an author-version comparison, arXiv:2410.04769v2, 15 October 2025, was separately acquired and bound [preprint]. Its physical pages 1 and 19–22 were read, and the title and decisive statement/proof pages were viewed. The public arXiv history checked during this review listed v2 as the latest version [history]. This is a bounded version check, not proof that no other correction exists.\par
\section*{Definitions and exact hypotheses}
Write \[R_\eta^D(\lambda)
 =\sum_{k\geq1}(\lambda-\lambda_k^D(\Omega))_+^\eta,
 \qquad
 L^{\mathrm{sc}}_{\eta,d}
 =\frac{\Gamma(1+\eta)}
 {(4\pi)^{d/2}\Gamma(1+\eta+d/2)}.\] For \(\eta=0\), the paper uses the strict counting function \(\#\{k:\lambda_k^D(\Omega)<\lambda\}\); the counterexample below has strictly positive exponents and avoids that convention entirely. The Laplacians are the quadratic-form realizations on an open set.\par
Equation reference: eq:hyp, eq:false, eq:printedf. Set \(a=d/2\). Lemma 3.3 assumes \(d\geq2\), \(\gamma>0\), an open set \(\Omega\subset\mathbb R^d\) of finite volume, \(c>0\), and \(0<\Lambda\leq\infty\), with \[
 R_\gamma^D(\lambda)
 \leq cL^{\mathrm{sc}}_{\gamma,d}|\Omega|\lambda^{\gamma+a}
 \quad(0\leq\lambda\leq\Lambda).\] It claims, for every \(q=\gamma'\in[0,\gamma]\), that \[
 R_q^D(\lambda)
 \leq c f_\gamma^D(q)L^{\mathrm{sc}}_{q,d}|\Omega|\lambda^{q+a},
 \qquad
 0\leq\lambda\leq\frac{q+a}{\gamma+a}\Lambda,\] where the printed coefficient is \[
 f_\gamma^D(q)=
 \frac{(\gamma+a)^{\gamma+a}}{(q+a)^{q+a}}
 \frac{\Gamma(1+q+a)\Gamma(1+\gamma)}
 {\Gamma(1+\gamma+a)\Gamma(1+q)}.\] There is no restriction \(\gamma\geq1\) or \(c\geq1\) in the statement.\par
\section*{A false extrapolation bound}
Problem 1 (P01: Major; Lemma 3.3).  The asserted coefficient in [eq:false] (\#eq:false) is too small for some allowed exponents. The exact source locator is printed p. 781, physical p. 21, Lemma 3.3. The omitted factor can be tracked through the minimization on printed p. 782, physical p. 22; the subsequent coefficient analysis continues on pp. 783–784. The acquired arXiv v2 has the same statement and minimization on physical pp. 19–20.\par
Equation reference: eq:parameters. Claim 1.  The square \(\Omega=(0,\pi)^2\), with \[
 d=2,\qquad \gamma=\frac14,\qquad q=\frac18,\qquad
 \Lambda=\frac52,\qquad c=\frac{2}{\pi\,5^{1/4}},\] satisfies every hypothesis of Lemma 3.3, but violates its conclusion at \(\lambda_*=9/4\).\par
Proof. The square is open, bounded, nonempty, convex and Lipschitz, with area \(\pi^2\). Separation of variables gives the Dirichlet eigenvalues \(k^2+\ell^2\), \(k,\ell\geq1\). The first eigenvalue is \(2\), simple, and the next is \(5\). Thus \(R_\gamma^D(\lambda)=0\) for \(0\leq\lambda\leq2\); for \(2<\lambda\leq5/2\) it is \((\lambda-2)^{1/4}\). Since \(L^{\mathrm{sc}}_{\gamma,2}|\Omega|=\pi/5\), the normalized ratio is \[F(\lambda)=
 \frac{5(\lambda-2)^{1/4}}{\pi\lambda^{5/4}},
 \qquad
 \frac{F'(\lambda)}{F(\lambda)}
 =\frac{5-2\lambda}{2\lambda(\lambda-2)}\geq0.\] Consequently its maximum over this entire energy interval is \(F(5/2)=2/(\pi5^{1/4})=c\). This proves [eq:hyp] (\#eq:hyp) for every required energy, including the endpoints; it is not merely a test at one energy.\par
The target is exactly the allowed endpoint: \[\frac{q+1}{\gamma+1}\Lambda
 =\frac{9/8}{5/4}\frac52=\frac94.\] Gamma-function cancellation in [eq:printedf] (\#eq:printedf) gives \[f_{1/4}^D(1/8)
 =\frac{(5/4)^{1/4}}{(9/8)^{1/8}}.\] Therefore \[\begin{aligned}
 R_{1/8}^D(9/4)&=(1/4)^{1/8}=2^{-1/4},\\
 c f_{1/4}^D(1/8)L^{\mathrm{sc}}_{1/8,2}|\Omega|(9/4)^{9/8}
 &=2^{-3/8}.\end{aligned}\] The quotient of the actual value by the claimed upper bound is \(2^{1/8}>1\). Both energies lie strictly between consecutive eigenvalues, so no counting-function or threshold convention can remove the contradiction. ◻\par
The exact identities, the logarithmic derivative, the permitted energy endpoint and the corrected equality were independently checked by a root-authored SymPy calculation. Its actual execution exited zero; the script, output, version information and execution receipt are retained with P01. The proof above is sufficient without a floating-point spectral calculation.\par
\section*{Repair and alternative readings}
Equation reference: eq:tau, eq:repair. The scalar comparison preceding the minimization is valid. For \(0<q<\gamma\), its normalized coefficient is \[
 \frac{q^q(\gamma-q)^{\gamma-q}}{\gamma^\gamma}
 \frac{(1+\tau)^{\gamma+a}}{\tau^{\gamma-q}}
 \frac{L^{\mathrm{sc}}_{\gamma,d}}{L^{\mathrm{sc}}_{q,d}},
 \qquad (1+\tau)\lambda\leq\Lambda.\] Indeed, maximizing \(x^q/(x+\delta)^\gamma\) for \(x\geq0\) gives \(x=q\delta/(\gamma-q)\) and the scalar coefficient used in Lemma 3.2. Minimizing [eq:tau] (\#eq:tau) gives \(\tau=(\gamma-q)/(q+a)\), exactly the paper’s value. Substitution, however, yields \[
 \widetilde f_\gamma^D(q)
 =\frac{q^q}{\gamma^\gamma}f_\gamma^D(q),\] with the same energy range as [eq:false] (\#eq:false). The factors \((\gamma-q)^{\gamma-q}\) cancel, but \(q^q/\gamma^\gamma\) does not. At the diagonal the original assumption applies; at \(q=0\) the corresponding limiting estimate uses \(0^0=1\). In the square example the extra factor is \(2^{1/8}\), so the repaired estimate is an equality.\par
This is not an alternative normalization of \(L^{\mathrm{sc}}\): the definition in (1.4) and its elementary \(d=2\) specialization were used explicitly. Nor can the extra factor always be discarded as a harmless improvement. It exceeds one in [eq:parameters] (\#eq:parameters). Restricting the lemma to \(\gamma=1\) would avoid this counterexample, but changes its stated range.\par
The correction does preserve the qualitative Dirichlet step in Proposition 3.1. Its hypotheses include \(\gamma_1>\gamma_0\geq0\), \(c<1\), \(\Lambda_1>\Lambda_0>0\), the semiclassical bound for every energy above \(\Lambda_0\) and every exponent in \([\gamma_0,\gamma_1]\), and a strict factor \(c\) at exponent \(\gamma_1\) below \(\Lambda_1\). By continuity of \(\widetilde f_{\gamma_1}^D(q)\) at \(q=\gamma_1\), where it equals one, choose \(q<\gamma_1\) sufficiently close to \(\gamma_1\) that \[c\widetilde f_{\gamma_1}^D(q)<1,\qquad
 \frac{q+a}{\gamma_1+a}\Lambda_1>\Lambda_0,\qquad q\geq\gamma_0.\] The repaired estimate covers the low energies and overlaps the assumed high-energy range. This proves the local exponent improvement used there, without the printed inverse-coefficient formula or any unproved monotonicity assertion for the new coefficient.\par
There is also direct counterevidence against a claim of failure of the main below-one inequality: for \(\gamma_1=1\) and \(0\leq q\leq1\), \(q^q\leq1\). Thus the printed coefficient is already a weaker valid bound in that special case. The proof of Theorems 1.1 and 1.2 on printed p. 786 explicitly starts at exponent one. The separate Neumann extrapolation in Lemma 3.4 does not have this omitted factor.\par
\section*{Dependency and impact boundaries}
The false general lemma is part of a core argument, not an isolated numerical illustration. The following dependencies are explicit in the acquired source:\par
1. Proposition 3.1 applies Lemma 3.3 on p. 784. Proposition 4.1 then invokes Proposition 3.1 on p. 786.\par
2. The implication (ii)\(\Rightarrow\)(i) in Theorem 4.5 uses Proposition 4.1 on p. 788. Theorem 4.5 characterizes the critical exponent and is an announced main ingredient of Theorem 1.5 on p. 799.\par
3. Section 5, fact (d), p. 791, cites Lemmas 3.3 and 3.4 for continuity of the optimal excess factors as the Riesz exponent varies.\par
The classification is Major because an asserted quantitative lemma is demonstrably false and its general-exponent derivation serves the paper’s central extrapolation and critical-exponent analysis. It does not assert whole-paper failure. The explicit repair above restores the local qualitative mechanism checked here; the report has not re-proved every downstream optimization or collapsing-domain theorem. A correction should replace the coefficient, revise the associated coefficient discussion and quantitative exponent choices, and check downstream uses.\par
Other checks did not yield a further supported substantive finding. The all-energy quantifier and the positive-exponent restriction in the example were checked directly. The product eigenvalue reduction in Lemma 5.3 and its use for Proposition 1.4 were followed, as were the strict exponent gap and the role of Theorem 1.8 in the Dirichlet compactness argument. The external two-term asymptotics and the full proof of Theorem 1.8 remain dependencies, not independently certified results of this review. Apparent local typographical slips were not promoted into additional Problems. No author program was run, and no mathematical defect is inferred from code availability.\par
\section*{Scope and references}
This initial report records one proved quantitative defect and its repair boundary. Scientific adjudication remains pending and no independent acceptance is claimed.\par
9 R. L. Frank and S. Larson, Semiclassical inequalities for Dirichlet and Neumann Laplacians on convex domains, Comm. Pure Appl. Math. 79 (2026), 762–822. doi:10.1002/cpa.70019 (https://doi.org/10.1002/cpa.70019). Acquired Chalmers version-of-record PDF. R. L. Frank and S. Larson, same title, arXiv:2410.04769v2 (https://arxiv.org/abs/2410.04769v2), 15 October 2025. Separately acquired 57-page author version; comparison of Lemma 3.3 on pp. 19–20. arXiv, Submission history for 2410.04769, https://arxiv.org/abs/2410.04769, accessed 29 October 2026. Chalmers publication metadata: https://research.chalmers.se/publication/548922.\par
\end{document}
